% !TeX program = xelatex
\documentclass[aspectratio=43,11pt,t]{beamer}
\usetheme{SYSU}
\usepackage{booktabs}

% 报告信息只在这里修改。
\title{学术报告题目}
\author{报告人姓名}
\institute{中山大学\quad 海洋工程与技术学院}
\date{\today}
\SYSUsetup{
  advisor={指导教师姓名},
  email={name@example.edu.cn},
  closing={请各位老师、同学批评指正！}
}

\begin{document}

\SYSUtitleframe
\SYSUoutlineframe

\section{研究背景}
\begin{frame}{Introduction}
  \begin{itemize}
    \item 研究问题与应用背景
      \begin{itemize}
        \item 说明研究对象、场景与基本假设
      \end{itemize}
    \item 现有方法与主要不足
    % 取消下一行的注释，即可按页逐步展示第三项。
    % \pause
    \item 本研究的目标与主要贡献
  \end{itemize}
\end{frame}

\section{研究方法}
\begin{frame}{Method}
  \SYSUheading{模型与符号}
  \[
    p(r,f)=\sum_{n=1}^{N} A_n(r,f)\,
      \exp\!\bigl(\mathrm{i}\,\phi_n(r,f)\bigr)
  \]
  \vspace{1mm}
  \begin{center}
    \begin{tabular}{lll}
      \toprule
      物理量 & 符号 & 单位 \\
      \midrule
      距离 & \(r\) & \(\mathrm{m}\) \\
      频率 & \(f\) & \(\mathrm{Hz}\) \\
      相位 & \(\phi_n\) & \(\mathrm{rad}\) \\
      \bottomrule
    \end{tabular}
  \end{center}
\end{frame}

\section{结果分析}
\begin{frame}{Analysis}
  \begin{columns}[T,onlytextwidth]
    \begin{column}{0.48\textwidth}
      \SYSUfigureplaceholder{结果图 A}
      \SYSUcaption{图 A：研究对象或基准结果}
    \end{column}
    \begin{column}{0.48\textwidth}
      \SYSUfigureplaceholder{结果图 B}
      \SYSUcaption{图 B：对比结果或误差分布}
    \end{column}
  \end{columns}
  \vspace{4mm}
  \SYSUheading{结果说明}
  说明图中变量、比较条件和主要发现。
\end{frame}

\section{总结展望}
\begin{frame}{Conclusion}
  \begin{enumerate}
    \item 主要结论
    \item 适用条件与当前局限
    \item 后续研究工作
  \end{enumerate}
\end{frame}

\SYSUthanksframe

\end{document}
